% ==============================================================================
% Sistema de pósits para tareas inconclusas y revisión personal (todonotes)
% ==============================================================================

% Para la entrega final sin pósits, comenta la primera línea y descomenta la segunda:
\usepackage[colorinlistoftodos, textsize=footnotesize]{todonotes}
% \usepackage[disable]{todonotes}

% --- Paleta diferenciada por tipo de anotación ---
\definecolor{positAmarilloBg}{HTML}{FFF3CD}
\definecolor{positAmarilloBorde}{HTML}{D39E00}

\definecolor{positAzulBg}{HTML}{DCEEFB}
\definecolor{positAzulBorde}{HTML}{2B6CB0}

\definecolor{positNaranjaBg}{HTML}{FFE8D6}
\definecolor{positNaranjaBorde}{HTML}{D96B27}

\definecolor{positVerdeBg}{HTML}{E2F3E8}
\definecolor{positVerdeBorde}{HTML}{2F855A}

% --- Comandos semánticos (sin caption={} para evitar conflictos en pgfkeys) ---
% Uso: nunca colocar \positrevision ni \positcita (al margen) dentro de table,
% figure o caption; en esos entornos pasar la opción [inline] o ubicarlos fuera.

\newcommand{\positinconcluso}[2][]{%
  \todo[inline, backgroundcolor=positAmarilloBg, bordercolor=positAmarilloBorde, linecolor=positAmarilloBorde, #1]{%
    \textbf{[INCONCLUSO]} #2%
  }%
}

\newcommand{\positrevision}[2][]{%
  \todo[backgroundcolor=positAzulBg, bordercolor=positAzulBorde, linecolor=positAzulBorde, #1]{%
    \textbf{[REVISIÓN]} #2%
  }%
}

\newcommand{\positdato}[2][]{%
  \todo[inline, backgroundcolor=positNaranjaBg, bordercolor=positNaranjaBorde, linecolor=positNaranjaBorde, #1]{%
    \textbf{[DATO PENDIENTE DE MEDIR]} #2%
  }%
}

\newcommand{\positcita}[2][]{%
  \todo[backgroundcolor=positVerdeBg, bordercolor=positVerdeBorde, linecolor=positVerdeBorde, #1]{%
    \textbf{[CITA APA]} #2%
  }%
}

\newcommand{\listadeposits}{%
  \listoftodos[Índice de pósits y tareas pendientes]%
}
